\chapter{Complex Analysis \& Applications}
\index{Complex Numbers}
\index{Euler's Formula}
\index{Argand Plane}
\index{De Moivre's Theorem}
\index{Phasors}
\index{Complex Impedance}
\index{Harmonic Waves}
\index{Cauchy-Riemann Equations}

% ==========================================================
% Chapter 4 Instructional Management Plan & Syllabus
% ==========================================================
\section*{Chapter 4 Instructional Management Plan}
\addcontentsline{toc}{section}{Chapter 4 Instructional Management Plan}

\begin{planbox}[colframe=rbruNavy, colbacktitle=rbruNavy]{1. Chapter Topics and Core Content}
\begin{enumerate}[topsep=2pt, itemsep=1pt, partopsep=0pt, parsep=0pt, leftmargin=1.5em]
    \item Algebraic and Geometric Foundations of Complex Numbers
    \item Euler's Formula, Polar Form, and Roots of Unity
    \item Physical Applications: Phasors and AC Impedance
    \item Worked Examples and Harmonic Wave Interference
    \item Summary of Key Formulas and Chapter Exercises
\end{enumerate}
\end{planbox}

\begin{planbox}[colframe=rbruGreen, colbacktitle=rbruGreen]{2. Behavioral Learning Objectives (OBE)}
\begin{enumerate}[topsep=2pt, itemsep=1pt, partopsep=0pt, parsep=0pt, leftmargin=1.5em]
    \item Represent complex numbers interchangeably in Cartesian, polar, and exponential forms.
    \item Apply de Moivre's theorem to compute roots of unity and trigonometric identities.
    \item Model AC circuits and damped oscillations using complex impedance and phasors.
    \item Verify analyticity of complex functions using Cauchy-Riemann equations.
\end{enumerate}
\end{planbox}

\newpage

\begin{planbox}[colframe=rbruNavy!85!black, colbacktitle=rbruNavy!85!black]{3. Teaching Methods and Instructional Activities}
\begin{enumerate}[topsep=2pt, itemsep=1pt, partopsep=0pt, parsep=0pt, leftmargin=1.5em]
    \item Lecture on the Argand plane, Euler's formula, and complex exponential mappings.
    \item Hands-on workshop deriving AC circuit impedance and phasor addition.
    \item Computational lab simulating phasor rotations and frequency responses in Python.
    \item Case study on resonance in RLC circuits and optical wave superposition.
\end{enumerate}
\end{planbox}

\begin{planbox}[colframe=rbruGold!90!black, colbacktitle=rbruGold!90!black]{4. Instructional Media and Educational Technology}
\begin{enumerate}[topsep=2pt, itemsep=1pt, partopsep=0pt, parsep=0pt, leftmargin=1.5em]
    \item Course textbook: Chapter 4 Complex Analysis \& Applications.
    \item Interactive digital slides displaying rotating phasor diagrams.
    \item Python scripts for Bode plots and complex impedance loci.
    \item Problem sets on complex algebra and AC circuit analysis.
\end{enumerate}
\end{planbox}

\begin{planbox}[colframe=rbruSlate, colbacktitle=rbruSlate]{5. Measurement and Learning Assessment}
\begin{enumerate}[topsep=2pt, itemsep=1pt, partopsep=0pt, parsep=0pt, leftmargin=1.5em]
    \item Evaluation of class participation and phasor diagram sketch exercises.
    \item Grading of homework problem sets on complex impedance and Euler's formula.
    \item Chapter 4 diagnostic quiz testing complex algebra and circuit analysis.
\end{enumerate}
\end{planbox}
\clearpage

\section{Algebraic and Geometric Foundations of Complex Numbers}
In classical mechanics, optics, and quantum mechanics, physical quantities that oscillate with specific phase relationships are most elegantly described using complex variables. A complex number $z \in \mathbb{C}$ is an ordered pair of real numbers $(x, y)$ written as:
\begin{equation}
z = x + iy, \quad i^2 = -1
\end{equation}
where $x = \text{Re}(z)$ is the real part, and $y = \text{Im}(z)$ is the imaginary part. The complex conjugate of $z$ is:
\begin{equation}
z^* = \bar{z} = x - iy
\end{equation}
The modulus (absolute magnitude) $|z|$ and argument $\theta = \text{arg}(z)$ are defined geometrically in the **Argand plane**:
\begin{equation}
|z| = \sqrt{z z^*} = \sqrt{x^2 + y^2}, \quad \theta = \arctan\left(\frac{y}{x}\right)
\end{equation}

\begin{maththeorem}[Euler's Formula and Polar Representation]
For any real angle $\theta \in \mathbb{R}$:
\begin{equation}
e^{i\theta} = \cos\theta + i\sin\theta
\end{equation}
Any complex number can thus be written in polar exponential form:
\begin{equation}
z = r e^{i\theta} = |z|(\cos\theta + i\sin\theta)
\end{equation}
When $\theta = \pi$, Euler's formula produces the renowned identity connecting fundamental constants:
\begin{equation}
e^{i\pi} + 1 = 0
\end{equation}
\end{maththeorem}

\subsection{De Moivre's Theorem and Roots of Unity}
For any integer $n \in \mathbb{Z}$:
\begin{equation}
(\cos\theta + i\sin\theta)^n = \cos(n\theta) + i\sin(n\theta) \iff (e^{i\theta})^n = e^{in\theta}
\end{equation}
The $n$ distinct $n$-th roots of a complex number $z_0 = r_0 e^{i\theta_0}$ are distributed evenly on a circle of radius $\sqrt[n]{r_0}$:
\begin{equation}
z_k = \sqrt[n]{r_0} \exp\left[i\left(\frac{\theta_0 + 2k\pi}{n}\right)\right], \quad k = 0, 1, \dots, n-1
\end{equation}

\section{Physical Applications: Phasors and AC Impedance}
In alternating current (AC) electrical networks and wave optics, sinusoidal voltages, currents, and fields of angular frequency $\omega$ are represented as the real part of complex rotating vectors (phasors):
\begin{equation}
V(t) = V_0 \cos(\omega t + \phi) = \text{Re}\left[\tilde{V}_0 e^{i\omega t}\right], \quad \tilde{V}_0 = V_0 e^{i\phi}
\end{equation}
Since time differentiation corresponds simply to multiplication by $i\omega$:
\begin{equation}
\frac{d}{dt}\left[\tilde{V}_0 e^{i\omega t}\right] = i\omega \left[\tilde{V}_0 e^{i\omega t}\right]
\end{equation}
differential equations governing linear dynamic circuits reduce to algebraic linear equations!

\begin{mathdefinition}[Complex Impedance]
The generalized complex impedance $\tilde{Z} \equiv \frac{\tilde{V}}{\tilde{I}}$ for fundamental circuit elements is:
\begin{itemize}
    \item Resistor ($R$): $\tilde{Z}_R = R$ (current and voltage are in phase).
    \item Inductor ($L$): $V = L\frac{dI}{dt} \implies \tilde{Z}_L = i\omega L = \omega L e^{i\pi/2}$ (voltage leads current by $90^\circ$).
    \item Capacitor ($C$): $I = C\frac{dV}{dt} \implies \tilde{Z}_C = \frac{1}{i\omega C} = -\frac{i}{\omega C} = \frac{1}{\omega C}e^{-i\pi/2}$ (voltage lags current by $90^\circ$).
\end{itemize}
\end{mathdefinition}

\begin{pitfallbox}[Taking Real Parts Before Nonlinear Operations]
The complex exponential trick works strictly for \textbf{linear} operators (differentiation, integration, linear superposition). When calculating nonlinear physical quantities such as power dissipation $P(t) = V(t)I(t)$ or electromagnetic energy density, one \textbf{must take the real parts first} before multiplying:
\begin{equation}
P(t) = \text{Re}[\tilde{V}(t)] \cdot \text{Re}[\tilde{I}(t)] \neq \text{Re}[\tilde{V}(t)\tilde{I}(t)]
\end{equation}
The time-averaged power for sinusoidal steady-state is $\langle P \rangle = \frac{1}{2}\text{Re}[\tilde{V}_0 \tilde{I}_0^*] = V_{\text{rms}} I_{\text{rms}} \cos\phi$.
\end{pitfallbox}

\section{Worked Examples and Physical Applications}

\begin{psctexample}[Series RLC Circuit Analysis via Complex Impedance]
A series $RLC$ circuit is driven by an alternating voltage source $V(t) = V_0 \cos(\omega t)$. Derive the analytical expression for the steady-state current amplitude $I_0$ and phase angle $\phi$ between voltage and current.

\textbf{\color{rbruNavy}Picture / Conceptualization:}
Resistor, inductor, and capacitor are wired in series. The driving voltage phasor is $\tilde{V} = V_0$.

\textbf{\color{rbruNavy}Strategy / Governing Laws:}
The total equivalent impedance of elements in series is the sum of their individual complex impedances:
\begin{equation}
\tilde{Z}_{\text{total}} = \tilde{Z}_R + \tilde{Z}_L + \tilde{Z}_C = R + i\left(\omega L - \frac{1}{\omega C}\right)
\end{equation}
The complex current phasor is given by Ohm's law in the frequency domain: $\tilde{I} = \frac{\tilde{V}}{\tilde{Z}_{\text{total}}}$.

\textbf{\color{rbruNavy}Calculation / Analytical Solution:}
1. Magnitude of total impedance:
\begin{equation}
|\tilde{Z}_{\text{total}}| = \sqrt{R^2 + \left(\omega L - \frac{1}{\omega C}\right)^2}
\end{equation}
2. Phase angle of impedance:
\begin{equation}
\phi = \arg(\tilde{Z}_{\text{total}}) = \arctan\left(\frac{\omega L - \frac{1}{\omega C}}{R}\right)
\end{equation}
3. Current phasor:
\begin{equation}
\tilde{I} = \frac{V_0}{|\tilde{Z}_{\text{total}}|e^{i\phi}} = \frac{V_0}{\sqrt{R^2 + \left(\omega L - \frac{1}{\omega C}\right)^2}}e^{-i\phi}
\end{equation}
4. Physical real current:
\begin{equation}
I(t) = \text{Re}[\tilde{I}e^{i\omega t}] = \frac{V_0}{\sqrt{R^2 + \left(\omega L - \frac{1}{\omega C}\right)^2}}\cos(\omega t - \phi)
\end{equation}

\textbf{\color{rbruNavy}Think / Physical Insights:}
At the resonance frequency $\omega_0 = \frac{1}{\sqrt{LC}}$, the inductive and capacitive reactances cancel exactly ($\omega L = 1/\omega C$). The impedance becomes purely real and minimal ($|\tilde{Z}| = R$), maximizing current amplitude $I_0 = V_0 / R$ with $\phi = 0$.
\qedexam
\end{psctexample}

\begin{pythonbox}[Complex Impedance and Resonance Response in Python]
import numpy as np
import matplotlib.pyplot as plt

# RLC Circuit Parameters
R = 50.0       # Ohms
L = 100e-3     # 100 mH
C = 10e-6      # 10 uF
V0 = 10.0      # Volts

omega_0 = 1.0 / np.sqrt(L * C)
freq = np.linspace(50, 400, 500)
omega = 2 * np.pi * freq

# Complex impedance calculation
Z = R + 1j * (omega * L - 1.0 / (omega * C))
I_mag = V0 / np.abs(Z)
phase_deg = np.angle(Z, deg=True)

print(f"Resonance Frequency f_0 = {omega_0 / (2*np.pi):.2f} Hz")
print(f"Max Current at resonance = {V0 / R:.3f} A")
\end{pythonbox}

\begin{psctexample}[Quantum Phase Interference and the Double-Slit Experiment]
A quantum particle passes simultaneously through two slits separated by distance $d$. The probability amplitudes from the two slits at an angle $\theta$ are complex numbers $\psi_1 = Ae^{i\phi_1}$ and $\psi_2 = Ae^{i\phi_2}$, where the phase difference is $\delta = \phi_2 - \phi_1 = \frac{2\pi d \sin\theta}{\lambda}$. Calculate (i) the total amplitude $\psi = \psi_1 + \psi_2$, (ii) the probability density $|\psi|^2$, and (iii) the positions of constructive and destructive interference.

\textbf{Solution:}

\textbf{(i) Total amplitude:}
\begin{equation}
\psi = Ae^{i\phi_1}(1 + e^{i\delta}) = Ae^{i\phi_1} \cdot 2\cos\!\left(\frac{\delta}{2}\right)e^{i\delta/2} = 2A\cos\!\left(\frac{\delta}{2}\right)e^{i(\phi_1 + \delta/2)}
\end{equation}

\textbf{(ii) Probability density (using modulus squared):}
\begin{equation}
\boxed{|\psi|^2 = |2A\cos(\delta/2)|^2 = 4A^2\cos^2\!\left(\frac{\pi d \sin\theta}{\lambda}\right)}
\end{equation}

\textbf{(iii) Interference conditions:}
\begin{align}
\text{Constructive: } \cos^2(\delta/2) = 1 \implies \delta = 2m\pi &\implies d\sin\theta = m\lambda, \quad m \in \mathbb{Z} \\
\text{Destructive: } \cos^2(\delta/2) = 0 \implies \delta = (2m+1)\pi &\implies d\sin\theta = \left(m + \tfrac{1}{2}\right)\lambda
\end{align}

\textbf{\color{rbruNavy}Think / Physical Insights:}
The interference pattern $\propto \cos^2(\delta/2)$ arises entirely from the complex algebra of probability amplitudes. Unlike classical probability (which would give $|\psi_1|^2 + |\psi_2|^2 = 2A^2$, a featureless sum), quantum mechanics predicts an oscillating pattern — an experimentally verified signature of wave-particle duality. The cross term $2\text{Re}[\psi_1\psi_2^*] = 2A^2\cos\delta$ is the quantum interference term.
\qedexam
\end{psctexample}

\section{Summary of Key Formulas}
\begin{table}[H]
\centering
\small
\begin{tabularx}{\textwidth}{llX}
\toprule
\textbf{Concept} & \textbf{Formula} & \textbf{Physical Application} \\
\midrule
Euler's Formula & $e^{i\theta} = \cos\theta + i\sin\theta$ & Wave kinematics, rotation operator in quantum mechanics \\
Complex Impedance & $\tilde{Z} = R + i(\omega L - \frac{1}{\omega C})$ & Frequency-domain analysis of AC filters and oscillators \\
Average Real Power & $\langle P \rangle = \frac{1}{2}\text{Re}[\tilde{V}_0 \tilde{I}_0^*]$ & Electrical power transmission, optical energy absorption \\
Cauchy-Riemann & $\frac{\partial u}{\partial x} = \frac{\partial v}{\partial y}, \; \frac{\partial u}{\partial y} = -\frac{\partial v}{\partial x}$ & 2D electrostatics, potential flow around aerofoils \\
Quantum Interference & $|\psi_1 + \psi_2|^2 = |\psi_1|^2 + |\psi_2|^2 + 2\text{Re}[\psi_1\psi_2^*]$ & Double-slit experiment, thin-film optical coatings \\
\bottomrule
\end{tabularx}
\caption{Summary of complex analysis relations in physical sciences.}
\label{tab:ch04_summary}
\end{table}

\section{Chapter Exercises}
\begin{enumerate}
    \item \textbf{Roots of Unity:} Determine all four roots of the equation $z^4 + 16 = 0$ in polar and Cartesian form, and plot them in the Argand plane.
    \item \textbf{Trigonometric Summation:} Using the geometric series formula with complex exponentials, evaluate the sum $\sum_{k=0}^n \cos(k\theta)$.
    \item \textbf{Parallel RLC Network:} Derive the equivalent complex admittance $\tilde{Y} = 1/\tilde{Z}$ for a parallel combination of $R$, $L$, and $C$, and find the anti-resonance condition.
    \item \textbf{Cauchy's Integral Formula:} Evaluate the contour integral $\oint_C \frac{e^{iz}}{z^2 + 1}\,dz$ where $C$ is the circle $|z| = 2$ traversed counterclockwise. Use the residue theorem and state which poles lie inside $C$.
    \item \textbf{Phasor Power Factor:} A series RL circuit has resistance $R = 50\;\Omega$ and inductance $L = 0.2\;\text{H}$. For a source voltage $\tilde{V} = 240e^{i0}$ V at $f = 50\;\text{Hz}$, calculate the complex power $\tilde{S} = \frac{1}{2}\tilde{V}\tilde{I}^*$ and the power factor $\cos\phi$.
\end{enumerate}
